\section{Grassmanians}

\bnote

    Subvarieties are assumed to be closed.

\enote

The focus of this discussion is on Schubert calculus, in particular the sort of problems it solves in
enumerative geometry.
In general the problem is to determine the number of spaces (of a given dimension) intersect some set
of given spaces.
We are dealing here with linear spaces and thus it is natural to consider a form of topology on linear
subspaces.

\bdefn

    Let $V$ be an $n$-dimensional vector space, then for $k\leq n$ we define $V$'s {\emphcolor Grassmanian}
    to be the space
    $$ G(k,V) = \set{\Lambda\leq V}[\hbox{$\Lambda$ is $k$-dimensional}] . $$

\edefn

The Grassmanian will of course be given the structure of a variety, in due time.
First of all, if $\Psi\colon V\cong V'$ is an isomorphism of vector spaces then their Grassmanians are in
bijection
$$ G(k,V) \lisoto G(k,V'),\qquad \Lambda\mapsto\Psi\Lambda , $$
and this is an isomorphism of varieties.
Writing
$$ G(k,n) = G(k,\bF^n) , $$
we see that for $n$-dimensional $V$, $G(k,V)\cong G(k,n)$.

Note further that there is an isomorphism
$$ G(k,V) \lisoto G(n-k,V^*),\qquad \Lambda\mapsto\ann\Lambda = \set{\phi\in V^*}[\phi\Lambda=0] , $$
which is again an isomorphism of varieties.
Thus we have $G(k,n)\cong G(n-k,n)$.

Now we proceed with endowing $G(k,V)$ a variety structure.
For $\Lambda\in G(k,V)$ consider the exterior power
$$ \bigwedge^k\Lambda \leq \bigwedge^kV , $$
which is $1$-dimensional.
Thus there is a natural map
$$ G(k,V) \to \bP\parens{\bigwedge^kV} \cong \bP^{\binom nk-1},\qquad \Lambda \mapsto \bigwedge^k\Lambda . $$
This map is called the {\emphcolor Pl\"ucker embedding}.
We will continue to use $N=\binom nk-1$.

\blemm

    The Pl\"ucker embedding is indeed injective, and its image in $\bP(V^k)$ is a subvariety.

\elemm

\bproof

    Choose basis vectors $v_1,\dots,v_k\in\Lambda$ so that $\eta=v_1\wedge\cdots\wedge v_k$ spans
    $\Lambda^k$.
    Considering the exterior algebra $\bigwedge V$, a vector $v\in V$ annihiltes $\eta$ (meaning
    $v\wedge\eta$) iff $v$ is in the span of $v_1,\dots,v_k$, i.e.\ $v\in\Lambda$.
    Thus we can determine $\Lambda$ from its image in $\bP(V^k)$ by
    $$ \Lambda = \set{v\in V}[v\wedge\eta=0] , $$
    so the Pl\"ucker embedding is injective.

    Extend a basis $v_1,\dots,v_k\in\Lambda$ to a basis of $V$.
    Then we can write a $\eta\in\bigwedge^kV$ as $\eta=\sum_I\eta_Iv_I$ for scalars $\eta_I$.
    Then

    Now consider the map
    $$ V \xtos{\wedge\eta} \bigwedge^{k+1}V , $$
    we claim that $\eta=v_1\wedge\cdots\wedge v_k$ for linearly independent $v_i\in V$ iff the kernel
    of this map has dimension at least $k$.
    If $\eta$ is of this form then $v_1,\dots,v_k$ are in the kernel.
    Otherwise suppose $v_1,\dots,v_k$ are linearly independent vectors in the kernel.
    We may extend these $v_i$ to a basis of $V$, so that there are scalars $\eta_I$ such that
    $$ \eta = \sum_I\eta_Iv_I,\qquad v_I = v_{i_1}\wedge\cdots\wedge v_{i_k}\hbox{ for $I=(i_1<\cdots<i_k)$} . $$
    Then we know
    $$ 0 = v_1\wedge\cdots\wedge v_k\wedge\eta = \sum_I\eta_Iv_1\wedge\cdots\wedge v_k\wedge v_I
    = \sum_{I\cap[k]=\emptyset}\eta_Iv_I , $$
    so that for any multi-index $I$ which is not $(1,\dots,k)$, we have $\eta_I=0$.
    Thus $\eta$ is a scalar multiple of $v_1\wedge\cdots\wedge v_k$, as desired.
    And thus the image of the Grassmanian is
    $$ G(k,V) = \set{\eta\in\bigwedge^kV}[\rank(V\xtos{\wedge\eta}\bigwedge^{k+1}V)\leq n-k] , $$
    which is an algebraic subset of $\bP(V^k)$ as it is the zeroset of all of the $(n-k+1)$st minors of
    the map $\bul\wedge\eta$.

    To show that $G(k,V)$ is irreducible, we note that there is a transitive action of $\gl(V)$ on
    $G(k,V)$.
    Choosing any $\Lambda\in G(k,V)$ we have that $\gl(V)\Lambda=G(k,V)$ is irreducible because
    it is the image of the irreducible $\gl(V)$ (it is an open subset of the irreducible $\bA^n$).
    \qed

\eproof

Let $e_1,\dots,e_n$ be a basis for $V$ and identify $V\cong\bF^n$ so that
$\Lambda$ is the row space of some matrix $k\times n$ matrix $A=(a_{ij})$.
Then $\bigwedge^kV$ is spanned by $e_I=e_{i_1}\wedge\cdots\wedge e_{i_k}$ for $I=(i_1<\cdots<i_k)$.
So for the basis $v_1,\dots,v_k\in\Lambda$ determined by $A$ there are scalars $p_I$ such that
$$ v_1\wedge\cdots\wedge v_k = \sum_Ip_Ie_I . $$
Then we see that $p_I$ is the determinant of the submatrix of $A$ whose columns are determined by $I$.
These are the {\emphcolor Pl\"ucker coordinates} of $A$.

\bexam

    Note that the Pl\"ucker embedding for $k=1$ gives an isomorphism $G(1,V)\cong\bP(V)$.
    So the most basic example of a Grassmanian is projective space.

\eexam

Because $G(k,V)$ is a variety it has an affine cover.
For an $(n-k)$-dimensional subspace $\Gamma\leq V$ define
$$ U_\Gamma = \set{\Lambda\in G}[\Lambda\cap\Gamma=0] $$
to be the set of all subspaces which do not intersect $\Gamma$.
This is a Zariski open subset of $G$, because taking $w_1,\dots,w_{n-k}$ a basis for $\Gamma$ and
setting $\eta=w_1\wedge\cdots\wedge w_{n-k}$, we see that
$$ U_\Gamma = \set{[\omega]\in G\subseteq\bP(V^k)}[\omega\wedge\eta\neq0] . $$
We now claim that $U_\Gamma$ is isomorphic to $\bA^{k(n-k)}$.
Let $\Omega\in U_\Gamma$ be a complementary $k$-space to $\Gamma$, so that $V=\Gamma\oplus\Omega$,
and for any $\Lambda\in U_\Gamma$ we have projections
$$ \Lambda \xtos{\pi_\Gamma} \Gamma,\qquad \Lambda \xtos{\pi_\Omega} \Omega $$
modulo $\Omega$ and $\Gamma$ respectively.
Because the kernel of $\pi_\Omega$ is $\Gamma$ we see that $\pi_\Omega$ is an isomorphism, and so
$\Lambda$ is the graph of
$$ \phi \colon \Omega \xtos{\pi_\Omega^{-1}} \Lambda \inj V \xtos{\pi_\Gamma} \Gamma . $$
Indeed its graph is
$$ \graph\phi = \set{\omega+\pi_\Gamma\pi_\Omega^{-1}\omega}[\omega\in\Omega] , $$
and any $v\in\Omega$ is in this graph via $\omega=\pi_\Omega(v)$, and they have the same dimension.

Conversely if $\phi\colon\Omega\to\Gamma$ is linear then its graph is complementary to $\Gamma$.
So we have a correspondence
$$ U_\Gamma \xtof{\Lambda\mapsto\pi_\Gamma\circ\pi_\Omega^{-1}}[\phi\mapsto\graph\phi] \hom(\Omega,\Gamma) . $$
And considering them in Pl\"ucker coordinates, we see that these are morphisms of varieties.
Because $\hom(\Omega,\Gamma)\cong\bA^{k(n-k)}$, we see that $G(k,V)$ is a variety of dimension $k(n-k)$.

\bprop

    The Grassmanian $G(k,V)$ is a projective variety of dimension $k(n-k)$.

\eprop

\bnote

    In S.\ L.\ Kleiman and Dan Laksov, they define the Grassmanian $G(k,n)$ as the collection of $k$-planes in
    $\bP^n$.
    This corresponds to our $G(k+1,n+1)$.
    We will denote this by
    $$ \bG(k,V) = \bG(k,\bP(V)) \cong G(k+1,n+1) . $$

\enote


